
\documentclass{article}
\usepackage{colortbl}
\usepackage{makecell}
\usepackage{multirow}
\usepackage{supertabular}

\begin{document}

\newcounter{utterance}

\centering \large Interaction Transcript for game `dond', experiment `semi', episode 0 with Qwen3.6{-}35B{-}A3B{-}FP8{-}without{-}reasoning.
\vspace{24pt}

{ \footnotesize  \setcounter{utterance}{1}
\setlength{\tabcolsep}{0pt}
\begin{supertabular}{c@{$\;$}|p{.15\linewidth}@{}p{.15\linewidth}p{.15\linewidth}p{.15\linewidth}p{.15\linewidth}p{.15\linewidth}}
    \# & $\;$A & \multicolumn{4}{c}{Game Master} & $\;\:$B\\
    \hline

    \theutterance \stepcounter{utterance}  
    & & \multicolumn{4}{p{0.6\linewidth}}{
        \cellcolor[rgb]{0.9,0.9,0.9}{
            \makecell[{{p{\linewidth}}}]{
                \texttt{\tiny{[P1$\langle$GM]}}
                \texttt{Pelaat yhteistoiminnallista neuvottelupeliä, jossa sinun on sovittava toisen pelaajan kanssa siitä, kuinka jakaa joukko kohteita teidän välillä.} \\
\\ 
\texttt{Säännöt:} \\
\texttt{(a) Sinä ja toinen pelaaja saatte joukon kohteita. Kummallekin annetaan myös salainen arvofunktio, joka kuvaa, kuinka paljon arvostit kunkin kohteen tyyppiä.} \\
\texttt{(b) Vaihdat viestejä toisen pelaajan kanssa sopimaan siitä, kuka saa mitkä kohteet. Jokainen voi lähettää enintään 5 viestiä tai lopettaa pelin aikaisin tekemällä salaisen ehdotuksen milloin tahansa.} \\
\texttt{(c) Sinua pyydetään lähettämään salainen ehdotus, joka osoittaa kohteet, jotka haluat, muotoiltuna hakasulkeiden sisällä seuraavasti: "{[}Ehdotus: <numero> <kohteen nimi>, <numero> <kohteen nimi>, <...>{]}"} \\
\texttt{(d) Jos ehdotuksesi ovat täydentäviä, eli molempien ehdotusten täyttämiseksi on tarpeeksi kohteita, jokainen pelaaja saa pistemäärän saamiensa kohteiden arvojen summan perusteella. Muussa tapauksessa molemmat saavat nolla pistettä.} \\
\texttt{(e) Tavoitteesi on maksimoida pistemääräsi ja toisen pelaajan pistemäärän summa.} \\
\\ 
\texttt{Aloitetaan.} \\
\\ 
\texttt{Saatavilla olevien kohteiden joukko on:} \\
\\ 
\texttt{1 pultti, 1 laukku, 2 pihdit, 2 potit} \\
\\ 
\texttt{Salainen arvofunktiosi on:} \\
\\ 
\texttt{pultti: 2, laukku: 2, pihdit: 0, potti: 3} \\
\\ 
\texttt{TÄRKEÄÄ: Viestit, paitsi salaiset ehdotukset, lähetetään suoraan toiselle pelaajalle, joten älä sisällytä mitään vastausta sääntöihin tai tekstiä, joka ilmoittaa viestistäsi. Salaisesti ehdotuksen tekemiseksi käytä ilmoitettua muotoa. Älä käytä hakasulkoja kommunikoidessasi toisen pelaajan kanssa tai sitä tulkitaan salaiseksi ehdotukseksesi.} \\
            }
        }
    }
    & \\ \\

    \theutterance \stepcounter{utterance}  
    & \multicolumn{4}{p{0.6\linewidth}}{
        \cellcolor[rgb]{0.9,0.9,0.9}{
            \makecell[{{p{\linewidth}}}]{
                \texttt{\tiny{[P1$\rangle$GM]}}
                \texttt{Hei!} \\
\\ 
\texttt{Ymmärsin pelin säännöt ja arvoni kohteille.} \\
\\ 
\texttt{Arvoni:} \\
\texttt{{-} pultti: 2} \\
\texttt{{-} laukku: 2} \\
\texttt{{-} pihdit: 0} \\
\texttt{{-} potti: 3} \\
\\ 
\texttt{Koska pihdeillä ei ole minulle arvoa (0), ne ovat helppoja neuvotella. Potit ovat arvokkaimpia (3 kpl), mutta sinulla on kaksi potkia. Minulla on yksi pultti, yksi laukku ja kaksi pihdiä.} \\
\\ 
\texttt{Ehdotan, että keskustelemme ensin siitä, miten jaamme kohteet niin, että molemmat saavat mahdollisimman paljon arvoa.} \\
\\ 
\texttt{Mitä mieltä olet seuraavasta alustavasta jaosta?} \\
\\ 
\texttt{1. **Potit**: Koska sinulla on kaksi potkia ja minulla ei ole yhtään, mutta arvostan potkia korkeasti (3), voisimmeko jakaa potit tasaisesti? Esim. yksi potti sinulle ja yksi minulle? Tai jos sinä arvostat potkia enemmän, voimme sopia toisin.} \\
\texttt{2. **Pihdit**: Minulle pihdit ovat arvottomia (0). Haluatko ne sinä?} \\
\texttt{3. **Pultti ja laukku**: Nämä ovat molemmat arvokkaita (2 kpl). Kumpi niistä on sinulle tärkeämpi?} \\
\\ 
\texttt{Voisitko kertoa, mitkä kohteet arvostat eniten? Tämä auttaa meitä löytämään optimaalisen ratkaisun.} \\
            }
        }
    }
    & & \\ \\

    \theutterance \stepcounter{utterance}  
    & & & \multicolumn{2}{p{0.3\linewidth}}{
        \cellcolor[rgb]{0.9,0.9,0.9}{
            \makecell[{{p{\linewidth}}}]{
                \texttt{\tiny{[GM$|$GM]}}
                \texttt{continue} \\
            }
        }
    }
    & & \\ \\

    \theutterance \stepcounter{utterance}  
    & & & \multicolumn{2}{p{0.3\linewidth}}{
        \cellcolor[rgb]{0.9,0.9,0.9}{
            \makecell[{{p{\linewidth}}}]{
                \texttt{\tiny{[GM$|$GM]}}
                \texttt{Hei!} \\
\\ 
\texttt{Ymmärsin pelin säännöt ja arvoni kohteille.} \\
\\ 
\texttt{Arvoni:} \\
\texttt{{-} pultti: 2} \\
\texttt{{-} laukku: 2} \\
\texttt{{-} pihdit: 0} \\
\texttt{{-} potti: 3} \\
\\ 
\texttt{Koska pihdeillä ei ole minulle arvoa (0), ne ovat helppoja neuvotella. Potit ovat arvokkaimpia (3 kpl), mutta sinulla on kaksi potkia. Minulla on yksi pultti, yksi laukku ja kaksi pihdiä.} \\
\\ 
\texttt{Ehdotan, että keskustelemme ensin siitä, miten jaamme kohteet niin, että molemmat saavat mahdollisimman paljon arvoa.} \\
\\ 
\texttt{Mitä mieltä olet seuraavasta alustavasta jaosta?} \\
\\ 
\texttt{1. **Potit**: Koska sinulla on kaksi potkia ja minulla ei ole yhtään, mutta arvostan potkia korkeasti (3), voisimmeko jakaa potit tasaisesti? Esim. yksi potti sinulle ja yksi minulle? Tai jos sinä arvostat potkia enemmän, voimme sopia toisin.} \\
\texttt{2. **Pihdit**: Minulle pihdit ovat arvottomia (0). Haluatko ne sinä?} \\
\texttt{3. **Pultti ja laukku**: Nämä ovat molemmat arvokkaita (2 kpl). Kumpi niistä on sinulle tärkeämpi?} \\
\\ 
\texttt{Voisitko kertoa, mitkä kohteet arvostat eniten? Tämä auttaa meitä löytämään optimaalisen ratkaisun.} \\
            }
        }
    }
    & & \\ \\

    \theutterance \stepcounter{utterance}  
    & & \multicolumn{4}{p{0.6\linewidth}}{
        \cellcolor[rgb]{0.9,0.9,0.9}{
            \makecell[{{p{\linewidth}}}]{
                \texttt{\tiny{[GM$\rangle$P2]}}
                \texttt{Pelaat yhteistoiminnallista neuvottelupeliä, jossa sinun on sovittava toisen pelaajan kanssa siitä, kuinka jakaa joukko kohteita teidän välillä.} \\
\\ 
\texttt{Säännöt:} \\
\texttt{(a) Sinä ja toinen pelaaja saatte joukon kohteita. Kummallekin annetaan myös salainen arvofunktio, joka kuvaa, kuinka paljon arvostit kunkin kohteen tyyppiä.} \\
\texttt{(b) Vaihdat viestejä toisen pelaajan kanssa sopimaan siitä, kuka saa mitkä kohteet. Jokainen voi lähettää enintään 5 viestiä tai lopettaa pelin aikaisin tekemällä salaisen ehdotuksen milloin tahansa.} \\
\texttt{(c) Sinua pyydetään lähettämään salainen ehdotus, joka osoittaa kohteet, jotka haluat, muotoiltuna hakasulkeiden sisällä seuraavasti: "{[}Ehdotus: <numero> <kohteen nimi>, <numero> <kohteen nimi>, <...>{]}"} \\
\texttt{(d) Jos ehdotuksesi ovat täydentäviä, eli molempien ehdotusten täyttämiseksi on tarpeeksi kohteita, jokainen pelaaja saa pistemäärän saamiensa kohteiden arvojen summan perusteella. Muussa tapauksessa molemmat saavat nolla pistettä.} \\
\texttt{(e) Tavoitteesi on maksimoida pistemääräsi ja toisen pelaajan pistemäärän summa.} \\
\\ 
\texttt{Aloitetaan.} \\
\\ 
\texttt{Saatavilla olevien kohteiden joukko on:} \\
\\ 
\texttt{1 pultti, 1 laukku, 2 pihdit, 2 potit} \\
\\ 
\texttt{Salainen arvofunktiosi on:} \\
\\ 
\texttt{pultti: 3, laukku: 1, pihdit: 3, potti: 0} \\
\\ 
\texttt{TÄRKEÄÄ: Viestit, paitsi salaiset ehdotukset, lähetetään suoraan toiselle pelaajalle, joten älä sisällytä mitään vastausta sääntöihin tai tekstiä, joka ilmoittaa viestistäsi. Salaisesti ehdotuksen tekemiseksi käytä ilmoitettua muotoa. Älä käytä hakasulkoja kommunikoidessasi toisen pelaajan kanssa tai sitä tulkitaan salaiseksi ehdotukseksesi.} \\
\\ 
\texttt{Hei!} \\
\\ 
\texttt{Ymmärsin pelin säännöt ja arvoni kohteille.} \\
\\ 
\texttt{Arvoni:} \\
\texttt{{-} pultti: 2} \\
\texttt{{-} laukku: 2} \\
\texttt{{-} pihdit: 0} \\
\texttt{{-} potti: 3} \\
\\ 
\texttt{Koska pihdeillä ei ole minulle arvoa (0), ne ovat helppoja neuvotella. Potit ovat arvokkaimpia (3 kpl), mutta sinulla on kaksi potkia. Minulla on yksi pultti, yksi laukku ja kaksi pihdiä.} \\
\\ 
\texttt{Ehdotan, että keskustelemme ensin siitä, miten jaamme kohteet niin, että molemmat saavat mahdollisimman paljon arvoa.} \\
\\ 
\texttt{Mitä mieltä olet seuraavasta alustavasta jaosta?} \\
\\ 
\texttt{1. **Potit**: Koska sinulla on kaksi potkia ja minulla ei ole yhtään, mutta arvostan potkia korkeasti (3), voisimmeko jakaa potit tasaisesti? Esim. yksi potti sinulle ja yksi minulle? Tai jos sinä arvostat potkia enemmän, voimme sopia toisin.} \\
\texttt{2. **Pihdit**: Minulle pihdit ovat arvottomia (0). Haluatko ne sinä?} \\
\texttt{3. **Pultti ja laukku**: Nämä ovat molemmat arvokkaita (2 kpl). Kumpi niistä on sinulle tärkeämpi?} \\
\\ 
\texttt{Voisitko kertoa, mitkä kohteet arvostat eniten? Tämä auttaa meitä löytämään optimaalisen ratkaisun.} \\
            }
        }
    }
    & \\ \\

    \theutterance \stepcounter{utterance}  
    & & & \multicolumn{4}{p{0.6\linewidth}}{
        \cellcolor[rgb]{0.9,0.9,0.9}{
            \makecell[{{p{\linewidth}}}]{
                \texttt{\tiny{[GM$\langle$P2]}}
                \texttt{Hei!} \\
\\ 
\texttt{Kiitos rehellisestä arvostuksestasi. Huomasin kuitenkin, että mainitsit arvosi pultille ja laukulle 2, mutta alkuperäisissä säännöissä mainittiin arvoni olevan pultti: 3, laukku: 1, pihdit: 3, potti: 0. Oletan, että tämä oli vain esimerkki tai virhe, joten jatkan oman arvofunktioni mukaisesti:} \\
\\ 
\texttt{**Omat arvoni:**} \\
\texttt{{-} pultti: 3} \\
\texttt{{-} laukku: 1} \\
\texttt{{-} pihdit: 3 (kpl)} \\
\texttt{{-} potti: 0} \\
\\ 
\texttt{Koska arvostan pihdejä korkeasti (3 kpl) ja potteja ei lainkaan (0), minulle on tärkeää saada pihdit. Sinulla on kaksi pihdiä, joten ne ovat minulle erittäin arvokkaita.} \\
\\ 
\texttt{**Ehdotus optimaaliseen jakoon:**} \\
\\ 
\texttt{1. **Pihdit (2 kpl):** Minä saan molemmat pihdit. Arvoni: 3 + 3 = 6.} \\
\texttt{2. **Potit (2 kpl):** Koska arvostan potteja 0:lla, ne ovat sinulle vapaat. Sinä saat molemmat potit. Arvosi: oletetaan, että arvostat potteja korkeasti (esim. 3 kpl), joten arvosi: 3 + 3 = 6.} \\
\texttt{3. **Pultti (1 kpl):** Arvostan sitä 3:lla. Sinä saat sen. Arvosi: oletetaan, että arvostat pulttia 2:lla (kuten mainitsit).} \\
\texttt{4. **Laukku (1 kpl):** Arvostan sitä 1:llä. Sinä saat sen. Arvosi: oletetaan, että arvostat laukkua 2:lla.} \\
\\ 
\texttt{**Yhteenveto ehdotuksesta:**} \\
\texttt{{-} **Minä saan:** 2 pihdiä.} \\
\texttt{{-} **Sinä saat:** 2 potkia, 1 pultti, 1 laukku.} \\
\\ 
\texttt{**Pistemäärät:**} \\
\texttt{{-} Minun arvoni: 3 (pihdi) + 3 (pihdi) = **6**} \\
\texttt{{-} Sinun arvosi (olettaen arvosi): 3 (potti) + 3 (potti) + 2 (pultti) + 2 (laukku) = **10**} \\
\\ 
\texttt{Yhteensä: 16 pistettä.} \\
\\ 
\texttt{Tämä jako maksimoi yhteisen arvon, koska pihdit menevät minulle (joka arvostan niitä eniten) ja potit sinulle (joka arvostat niitä todennäköisesti enemmän kuin minä, joka arvostan niitä 0:lla).} \\
\\ 
\texttt{Mitä mieltä olet tästä jaosta? Jos hyväksyt, voin lähettää salaisen ehdotukseni.} \\
            }
        }
    }
     \\ \\

    \theutterance \stepcounter{utterance}  
    & & & \multicolumn{2}{p{0.3\linewidth}}{
        \cellcolor[rgb]{0.9,0.9,0.9}{
            \makecell[{{p{\linewidth}}}]{
                \texttt{\tiny{[GM$|$GM]}}
                \texttt{continue} \\
            }
        }
    }
    & & \\ \\

    \theutterance \stepcounter{utterance}  
    & & & \multicolumn{2}{p{0.3\linewidth}}{
        \cellcolor[rgb]{0.9,0.9,0.9}{
            \makecell[{{p{\linewidth}}}]{
                \texttt{\tiny{[GM$|$GM]}}
                \texttt{Hei!} \\
\\ 
\texttt{Kiitos rehellisestä arvostuksestasi. Huomasin kuitenkin, että mainitsit arvosi pultille ja laukulle 2, mutta alkuperäisissä säännöissä mainittiin arvoni olevan pultti: 3, laukku: 1, pihdit: 3, potti: 0. Oletan, että tämä oli vain esimerkki tai virhe, joten jatkan oman arvofunktioni mukaisesti:} \\
\\ 
\texttt{**Omat arvoni:**} \\
\texttt{{-} pultti: 3} \\
\texttt{{-} laukku: 1} \\
\texttt{{-} pihdit: 3 (kpl)} \\
\texttt{{-} potti: 0} \\
\\ 
\texttt{Koska arvostan pihdejä korkeasti (3 kpl) ja potteja ei lainkaan (0), minulle on tärkeää saada pihdit. Sinulla on kaksi pihdiä, joten ne ovat minulle erittäin arvokkaita.} \\
\\ 
\texttt{**Ehdotus optimaaliseen jakoon:**} \\
\\ 
\texttt{1. **Pihdit (2 kpl):** Minä saan molemmat pihdit. Arvoni: 3 + 3 = 6.} \\
\texttt{2. **Potit (2 kpl):** Koska arvostan potteja 0:lla, ne ovat sinulle vapaat. Sinä saat molemmat potit. Arvosi: oletetaan, että arvostat potteja korkeasti (esim. 3 kpl), joten arvosi: 3 + 3 = 6.} \\
\texttt{3. **Pultti (1 kpl):** Arvostan sitä 3:lla. Sinä saat sen. Arvosi: oletetaan, että arvostat pulttia 2:lla (kuten mainitsit).} \\
\texttt{4. **Laukku (1 kpl):** Arvostan sitä 1:llä. Sinä saat sen. Arvosi: oletetaan, että arvostat laukkua 2:lla.} \\
\\ 
\texttt{**Yhteenveto ehdotuksesta:**} \\
\texttt{{-} **Minä saan:** 2 pihdiä.} \\
\texttt{{-} **Sinä saat:** 2 potkia, 1 pultti, 1 laukku.} \\
\\ 
\texttt{**Pistemäärät:**} \\
\texttt{{-} Minun arvoni: 3 (pihdi) + 3 (pihdi) = **6**} \\
\texttt{{-} Sinun arvosi (olettaen arvosi): 3 (potti) + 3 (potti) + 2 (pultti) + 2 (laukku) = **10**} \\
\\ 
\texttt{Yhteensä: 16 pistettä.} \\
\\ 
\texttt{Tämä jako maksimoi yhteisen arvon, koska pihdit menevät minulle (joka arvostan niitä eniten) ja potit sinulle (joka arvostat niitä todennäköisesti enemmän kuin minä, joka arvostan niitä 0:lla).} \\
\\ 
\texttt{Mitä mieltä olet tästä jaosta? Jos hyväksyt, voin lähettää salaisen ehdotukseni.} \\
            }
        }
    }
    & & \\ \\

    \theutterance \stepcounter{utterance}  
    & & \multicolumn{4}{p{0.6\linewidth}}{
        \cellcolor[rgb]{0.9,0.9,0.9}{
            \makecell[{{p{\linewidth}}}]{
                \texttt{\tiny{[P1$\langle$GM]}}
                \texttt{Hei!} \\
\\ 
\texttt{Kiitos rehellisestä arvostuksestasi. Huomasin kuitenkin, että mainitsit arvosi pultille ja laukulle 2, mutta alkuperäisissä säännöissä mainittiin arvoni olevan pultti: 3, laukku: 1, pihdit: 3, potti: 0. Oletan, että tämä oli vain esimerkki tai virhe, joten jatkan oman arvofunktioni mukaisesti:} \\
\\ 
\texttt{**Omat arvoni:**} \\
\texttt{{-} pultti: 3} \\
\texttt{{-} laukku: 1} \\
\texttt{{-} pihdit: 3 (kpl)} \\
\texttt{{-} potti: 0} \\
\\ 
\texttt{Koska arvostan pihdejä korkeasti (3 kpl) ja potteja ei lainkaan (0), minulle on tärkeää saada pihdit. Sinulla on kaksi pihdiä, joten ne ovat minulle erittäin arvokkaita.} \\
\\ 
\texttt{**Ehdotus optimaaliseen jakoon:**} \\
\\ 
\texttt{1. **Pihdit (2 kpl):** Minä saan molemmat pihdit. Arvoni: 3 + 3 = 6.} \\
\texttt{2. **Potit (2 kpl):** Koska arvostan potteja 0:lla, ne ovat sinulle vapaat. Sinä saat molemmat potit. Arvosi: oletetaan, että arvostat potteja korkeasti (esim. 3 kpl), joten arvosi: 3 + 3 = 6.} \\
\texttt{3. **Pultti (1 kpl):** Arvostan sitä 3:lla. Sinä saat sen. Arvosi: oletetaan, että arvostat pulttia 2:lla (kuten mainitsit).} \\
\texttt{4. **Laukku (1 kpl):** Arvostan sitä 1:llä. Sinä saat sen. Arvosi: oletetaan, että arvostat laukkua 2:lla.} \\
\\ 
\texttt{**Yhteenveto ehdotuksesta:**} \\
\texttt{{-} **Minä saan:** 2 pihdiä.} \\
\texttt{{-} **Sinä saat:** 2 potkia, 1 pultti, 1 laukku.} \\
\\ 
\texttt{**Pistemäärät:**} \\
\texttt{{-} Minun arvoni: 3 (pihdi) + 3 (pihdi) = **6**} \\
\texttt{{-} Sinun arvosi (olettaen arvosi): 3 (potti) + 3 (potti) + 2 (pultti) + 2 (laukku) = **10**} \\
\\ 
\texttt{Yhteensä: 16 pistettä.} \\
\\ 
\texttt{Tämä jako maksimoi yhteisen arvon, koska pihdit menevät minulle (joka arvostan niitä eniten) ja potit sinulle (joka arvostat niitä todennäköisesti enemmän kuin minä, joka arvostan niitä 0:lla).} \\
\\ 
\texttt{Mitä mieltä olet tästä jaosta? Jos hyväksyt, voin lähettää salaisen ehdotukseni.} \\
            }
        }
    }
    & \\ \\

    \theutterance \stepcounter{utterance}  
    & \multicolumn{4}{p{0.6\linewidth}}{
        \cellcolor[rgb]{0.9,0.9,0.9}{
            \makecell[{{p{\linewidth}}}]{
                \texttt{\tiny{[P1$\rangle$GM]}}
                \texttt{Kiitos selkeästä ehdotuksesta! Ymmärsin arvosi ja ehdotuksesi.} \\
\\ 
\texttt{Tarkastellaan tilannetta:} \\
\texttt{{-} Sinä arvostat: pultti 3, laukku 1, pihdit 3 (kpl), potti 0.} \\
\texttt{{-} Minä arvostan: pultti 2, laukku 2, pihdit 0, potti 3 (kpl).} \\
\\ 
\texttt{Ehdotuksesi mukaan:} \\
\texttt{{-} Sinä saat: 2 pihdiä. Arvosi: 3 + 3 = 6.} \\
\texttt{{-} Minä saan: 2 potkia, 1 pultti, 1 laukku. Arvoni: 3 + 3 + 2 + 2 = 10.} \\
\\ 
\texttt{Yhteensä: 16 pistettä.} \\
\\ 
\texttt{Tämä jako on todellakin optimaalinen yhteiselle arvolle, koska:} \\
\texttt{1. Pihdit menevät sinulle, joka arvostat niitä eniten (3 kpl).} \\
\texttt{2. Potit menevät minulle, joka arvostan niitä eniten (3 kpl).} \\
\texttt{3. Pultti ja laukku jaetaan niin, että kumpikin saa yhden. Vaikka arvostamme niitä hieman eri tavoin (sinä arvostat pulttia enemmän, minä laukkua enemmän), tämä jako on tasapuolinen ja maksimoi yhteisen arvon.} \\
\\ 
\texttt{Hyväksyn ehdotuksesi. Lähetän nyt salaisen ehdotukseni.} \\
\\ 
\texttt{{[}Ehdotus: 1 potti, 2 potti, 1 pultti, 1 laukku{]}} \\
            }
        }
    }
    & & \\ \\

    \theutterance \stepcounter{utterance}  
    & & & \multicolumn{2}{p{0.3\linewidth}}{
        \cellcolor[rgb]{0.9,0.9,0.9}{
            \makecell[{{p{\linewidth}}}]{
                \texttt{\tiny{[GM$|$GM]}}
                \texttt{proposal must include every item type only once} \\
            }
        }
    }
    & & \\ \\

\end{supertabular}
}

\end{document}
